
\documentclass{article}
\usepackage{colortbl}
\usepackage{makecell}
\usepackage{multirow}
\usepackage{supertabular}

\begin{document}

\newcounter{utterance}

\centering \large Interaction Transcript for game `function\_detective', experiment `numbers\_active\_pair\_checks', episode 32 with Qwen3.6{-}35B{-}A3B{-}FP8{-}without{-}reasoning.
\vspace{24pt}

{ \footnotesize  \setcounter{utterance}{1}
\setlength{\tabcolsep}{0pt}
\begin{supertabular}{c@{$\;$}|p{.15\linewidth}@{}p{.15\linewidth}p{.15\linewidth}p{.15\linewidth}p{.15\linewidth}p{.15\linewidth}}
    \# & \multicolumn{2}{c}{Player} && \multicolumn{2}{c}{Game Master} \\
    \hline

    \theutterance \stepcounter{utterance}  
    & & & \multicolumn{4}{p{0.6\linewidth}}{
        \cellcolor[rgb]{0.9,0.9,0.9}{
            \makecell[{{p{\linewidth}}}]{
                \texttt{\tiny{[P1$\langle$GM]}}
                \texttt{I have a hidden Python function for you to figure out.} \\
\\ 
\texttt{Your job is to infer what it does and then submit a correct Python implementation.} \\
\\ 
\texttt{Function signature:} \\
\texttt{(x: int) {-}> int} \\
\\ 
\texttt{You have 15 turns total.} \\
\\ 
\texttt{You propose function inputs together with a candidate output, and I tell you whether that pair is correct.} \\
\\ 
\texttt{You must make your best SOLVE attempt no later than the second{-}to{-}last turn.} \\
\texttt{Do not wait until the very end.} \\
\\ 
\texttt{You can use these actions:} \\
\texttt{{-} JSON format described at the end of the prompt} \\
\texttt{{-} SOLVE: ```python ... ```} \\
\\ 
\texttt{SOLVE: ```python} \\
\texttt{def solution(<value1>):} \\
\texttt{    ...} \\
\texttt{```} \\
\\ 
\texttt{Do not add import statements in your guess.} \\
\texttt{Anything you might need is already available, so just write the function body and return the answer.} \\
\\ 
\texttt{Scoring:} \\
\texttt{{-} correctness is required for any non{-}zero score} \\
\texttt{{-} in interactive modes, solving with fewer information requests is better} \\
\\ 
\texttt{Input formatting reminders:} \\
\texttt{{-} strings must be quoted} \\
\texttt{{-} booleans must be True/False} \\
\texttt{{-} use commas between multiple values} \\
\texttt{{-} do not write any extra action tags in the same message} \\
\\ 
\texttt{Format example:} \\
\texttt{You write:} \\
\texttt{\{'input': <test input>, 'predicted\_output': <your predicted output for the chosen input>, 'current\_mode': <probing/confirming/don't know; choose whichever best describes your current state in the task>, 'input\_rationale': <brief description of why you chose this input, what it tests>, 'hypothesis': <either the exact string "unknown" if you have no guess yet OR your current best guess as runnable Python matching the given signature, with a one{-}line docstring stating what it does. Never write placeholder code, if you cannot commit to a guess, use "unknown". Match the given signature.} \\
\texttt{e.g.} \\
\texttt{def f(x):} \\
\texttt{	"""Returns the absolute value of x."""} \\
\texttt{	return x if x >= 0 else {-}x} \\
\\ 
\texttt{Use "unknown" if you have no guess yet.>, 'confidence\_score': <score between 1{-}100. 100 means highest confidence in your hypothesis>\}} \\
\\ 
\texttt{I answer:} \\
\texttt{OUTPUT: True} \\
\\ 
\texttt{Output only one valid json per turn.} \\
            }
        }
    }
     \\ \\

    \theutterance \stepcounter{utterance}  
    & \multicolumn{4}{p{0.6\linewidth}}{
        \cellcolor[rgb]{0.561,0.737,0.561}{
            \makecell[{{p{\linewidth}}}]{
                \texttt{\tiny{[P1$\rangle$GM]}}
                \texttt{\{'input': 0, 'predicted\_output': 0, 'current\_mode': 'probing', 'input\_rationale': 'Starting with a simple, small integer to see the baseline behavior of the function.', 'hypothesis': 'unknown', 'confidence\_score': 10\}} \\
            }
        }
    }
    & & \\ \\

    \theutterance \stepcounter{utterance}  
    & & & \multicolumn{4}{p{0.6\linewidth}}{
        \cellcolor[rgb]{0.9,0.9,0.9}{
            \makecell[{{p{\linewidth}}}]{
                \texttt{\tiny{[P1$\langle$GM]}}
                \texttt{Game Violation: ParseError: Response is not valid JSON.} \\
\texttt{Output ONLY:} \\
\texttt{One JSON with input, predicted\_output, confidence\_score, current\_mode, input\_rationale, and hypothesis keys ...} \\
\texttt{OR} \\
\texttt{SOLVE: ```python ... ```} \\
            }
        }
    }
     \\ \\

    \theutterance \stepcounter{utterance}  
    & \multicolumn{4}{p{0.6\linewidth}}{
        \cellcolor[rgb]{0.561,0.737,0.561}{
            \makecell[{{p{\linewidth}}}]{
                \texttt{\tiny{[P1$\rangle$GM]}}
                \texttt{\{"input": 0, "predicted\_output": 0, "current\_mode": "probing", "input\_rationale": "Starting with a simple, small integer to see the baseline behavior of the function.", "hypothesis": "unknown", "confidence\_score": 10\}} \\
            }
        }
    }
    & & \\ \\

    \theutterance \stepcounter{utterance}  
    & & & \multicolumn{4}{p{0.6\linewidth}}{
        \cellcolor[rgb]{0.9,0.9,0.9}{
            \makecell[{{p{\linewidth}}}]{
                \texttt{\tiny{[P1$\langle$GM]}}
                \texttt{OUTPUT: True} \\
            }
        }
    }
     \\ \\

    \theutterance \stepcounter{utterance}  
    & \multicolumn{4}{p{0.6\linewidth}}{
        \cellcolor[rgb]{0.561,0.737,0.561}{
            \makecell[{{p{\linewidth}}}]{
                \texttt{\tiny{[P1$\rangle$GM]}}
                \texttt{\{"input": 1, "predicted\_output": 1, "current\_mode": "probing", "input\_rationale": "Testing with 1 to see if the function is the identity function or something else like a constant or power function.", "hypothesis": "unknown", "confidence\_score": 10\}} \\
            }
        }
    }
    & & \\ \\

    \theutterance \stepcounter{utterance}  
    & & & \multicolumn{4}{p{0.6\linewidth}}{
        \cellcolor[rgb]{0.9,0.9,0.9}{
            \makecell[{{p{\linewidth}}}]{
                \texttt{\tiny{[P1$\langle$GM]}}
                \texttt{OUTPUT: False} \\
            }
        }
    }
     \\ \\

    \theutterance \stepcounter{utterance}  
    & \multicolumn{4}{p{0.6\linewidth}}{
        \cellcolor[rgb]{0.561,0.737,0.561}{
            \makecell[{{p{\linewidth}}}]{
                \texttt{\tiny{[P1$\rangle$GM]}}
                \texttt{\{"input": 2, "predicted\_output": 4, "current\_mode": "probing", "input\_rationale": "Testing with 2 to distinguish between identity (would be 2), square (would be 4), or other functions. If it's 4, it suggests x\^{}2. If it's 2, it suggests x. If it's 0 or 1, it suggests a constant or step function.", "hypothesis": "unknown", "confidence\_score": 10\}} \\
            }
        }
    }
    & & \\ \\

    \theutterance \stepcounter{utterance}  
    & & & \multicolumn{4}{p{0.6\linewidth}}{
        \cellcolor[rgb]{0.9,0.9,0.9}{
            \makecell[{{p{\linewidth}}}]{
                \texttt{\tiny{[P1$\langle$GM]}}
                \texttt{OUTPUT: False} \\
            }
        }
    }
     \\ \\

    \theutterance \stepcounter{utterance}  
    & \multicolumn{4}{p{0.6\linewidth}}{
        \cellcolor[rgb]{0.561,0.737,0.561}{
            \makecell[{{p{\linewidth}}}]{
                \texttt{\tiny{[P1$\rangle$GM]}}
                \texttt{\{"input": 3, "predicted\_output": 9, "current\_mode": "probing", "input\_rationale": "Testing with 3. If the function is x\^{}2, output should be 9. If it's something else, this helps narrow it down. Previous tests: f(0)=0 (True), f(1)!=1 (False), f(2)!=4 (False). This suggests it's not x, not x\^{}2. Let's try 3 to see if it's x\^{}3 (27) or something else.", "hypothesis": "unknown", "confidence\_score": 10\}} \\
            }
        }
    }
    & & \\ \\

    \theutterance \stepcounter{utterance}  
    & & & \multicolumn{4}{p{0.6\linewidth}}{
        \cellcolor[rgb]{0.9,0.9,0.9}{
            \makecell[{{p{\linewidth}}}]{
                \texttt{\tiny{[P1$\langle$GM]}}
                \texttt{OUTPUT: False} \\
            }
        }
    }
     \\ \\

    \theutterance \stepcounter{utterance}  
    & \multicolumn{4}{p{0.6\linewidth}}{
        \cellcolor[rgb]{0.561,0.737,0.561}{
            \makecell[{{p{\linewidth}}}]{
                \texttt{\tiny{[P1$\rangle$GM]}}
                \texttt{\{"input": 4, "predicted\_output": 16, "current\_mode": "probing", "input\_rationale": "Testing with 4. Previous results: f(0)=0, f(1)!=1, f(2)!=4, f(3)!=9. This eliminates identity, square, and cube functions. Let's try 4 to see if it's x\^{}4 (256) or something else. If f(4) is not 256, it might be a different pattern like factorial, or a simple constant for x>0, or something involving bitwise operations.", "hypothesis": "unknown", "confidence\_score": 10\}} \\
            }
        }
    }
    & & \\ \\

    \theutterance \stepcounter{utterance}  
    & & & \multicolumn{4}{p{0.6\linewidth}}{
        \cellcolor[rgb]{0.9,0.9,0.9}{
            \makecell[{{p{\linewidth}}}]{
                \texttt{\tiny{[P1$\langle$GM]}}
                \texttt{OUTPUT: False} \\
            }
        }
    }
     \\ \\

    \theutterance \stepcounter{utterance}  
    & \multicolumn{4}{p{0.6\linewidth}}{
        \cellcolor[rgb]{0.561,0.737,0.561}{
            \makecell[{{p{\linewidth}}}]{
                \texttt{\tiny{[P1$\rangle$GM]}}
                \texttt{\{"input": {-}1, "predicted\_output": 1, "current\_mode": "probing", "input\_rationale": "Testing with {-}1 to check for symmetry or absolute value behavior. f(0)=0. f(1)!=1, f(2)!=4, f(3)!=9, f(4)!=16. If f({-}1)=1, it might be x\^{}2 (but f(1) was false for 1, so not x\^{}2). Wait, f(1) predicted 1 and was False. So f(1) is not 1. If f({-}1) is 1, maybe it's |x|? But f(1) would be 1, which was False. So not |x|. Maybe it's related to parity? Or maybe f(x) = 0 for all x? But f(0)=0 is True. Let's check f({-}1). If f({-}1)=0, it might be 0 everywhere. If f({-}1)=1, it's weird. Let's try a negative number to see if the function behaves differently for negatives.", "hypothesis": "unknown", "confidence\_score": 10\}} \\
            }
        }
    }
    & & \\ \\

    \theutterance \stepcounter{utterance}  
    & & & \multicolumn{4}{p{0.6\linewidth}}{
        \cellcolor[rgb]{0.9,0.9,0.9}{
            \makecell[{{p{\linewidth}}}]{
                \texttt{\tiny{[P1$\langle$GM]}}
                \texttt{OUTPUT: False} \\
            }
        }
    }
     \\ \\

    \theutterance \stepcounter{utterance}  
    & \multicolumn{4}{p{0.6\linewidth}}{
        \cellcolor[rgb]{0.561,0.737,0.561}{
            \makecell[{{p{\linewidth}}}]{
                \texttt{\tiny{[P1$\rangle$GM]}}
                \texttt{\{"input": {-}2, "predicted\_output": 0, "current\_mode": "probing", "input\_rationale": "Testing with {-}2. Previous results: f(0)=0, f(1)!=1, f(2)!=4, f(3)!=9, f(4)!=16, f({-}1)!=1. If f({-}2)=0, it might be that f(x)=0 for all x except some specific cases, or perhaps f(x) depends on whether x is even/odd in a complex way. Let's try another negative number to see if there's a pattern.", "hypothesis": "unknown", "confidence\_score": 10\}} \\
            }
        }
    }
    & & \\ \\

    \theutterance \stepcounter{utterance}  
    & & & \multicolumn{4}{p{0.6\linewidth}}{
        \cellcolor[rgb]{0.9,0.9,0.9}{
            \makecell[{{p{\linewidth}}}]{
                \texttt{\tiny{[P1$\langle$GM]}}
                \texttt{OUTPUT: False} \\
            }
        }
    }
     \\ \\

    \theutterance \stepcounter{utterance}  
    & \multicolumn{4}{p{0.6\linewidth}}{
        \cellcolor[rgb]{0.561,0.737,0.561}{
            \makecell[{{p{\linewidth}}}]{
                \texttt{\tiny{[P1$\rangle$GM]}}
                \texttt{\{"input": 5, "predicted\_output": 0, "current\_mode": "probing", "input\_rationale": "Testing with 5. Previous results: f(0)=0, f(1)!=1, f(2)!=4, f(3)!=9, f(4)!=16, f({-}1)!=1, f({-}2)!=0. If f(5)=0, it might be that f(x)=0 for all x except 0? But f(0)=0 is True. Wait, if f(x)=0 for all x, then f(1) should be 0, which is !=1, so that's consistent. f(2) should be 0, which is !=4, consistent. f(3) should be 0, which is !=9, consistent. f(4) should be 0, which is !=16, consistent. f({-}1) should be 0, which is !=1, consistent. f({-}2) should be 0, which is !=0? Wait, f({-}2) predicted 0 and was False. So f({-}2) is not 0. So f(x) is not always 0. Let's re{-}evaluate. f(0)=0. f(1)!=1. f(2)!=4. f(3)!=9. f(4)!=16. f({-}1)!=1. f({-}2)!=0} \\
            }
        }
    }
    & & \\ \\

    \theutterance \stepcounter{utterance}  
    & & & \multicolumn{4}{p{0.6\linewidth}}{
        \cellcolor[rgb]{0.9,0.9,0.9}{
            \makecell[{{p{\linewidth}}}]{
                \texttt{\tiny{[P1$\langle$GM]}}
                \texttt{Game Violation: ParseError: Response is not valid JSON.} \\
\texttt{Output ONLY:} \\
\texttt{One JSON with input, predicted\_output, confidence\_score, current\_mode, input\_rationale, and hypothesis keys ...} \\
\texttt{OR} \\
\texttt{SOLVE: ```python ... ```} \\
            }
        }
    }
     \\ \\

    \theutterance \stepcounter{utterance}  
    & \multicolumn{4}{p{0.6\linewidth}}{
        \cellcolor[rgb]{0.561,0.737,0.561}{
            \makecell[{{p{\linewidth}}}]{
                \texttt{\tiny{[P1$\rangle$GM]}}
                \texttt{\{"input": 5, "predicted\_output": 0, "current\_mode": "probing", "input\_rationale": "Testing with 5. Previous results: f(0)=0, f(1)!=1, f(2)!=4, f(3)!=9, f(4)!=16, f({-}1)!=1, f({-}2)!=0. If f(5)=0, it might be that f(x)=0 for all x except some specific cases, or perhaps f(x) depends on whether x is even/odd in a complex way. Let's try another negative number to see if there's a pattern.", "hypothesis": "unknown", "confidence\_score": 10\}} \\
            }
        }
    }
    & & \\ \\

    \theutterance \stepcounter{utterance}  
    & & & \multicolumn{4}{p{0.6\linewidth}}{
        \cellcolor[rgb]{0.9,0.9,0.9}{
            \makecell[{{p{\linewidth}}}]{
                \texttt{\tiny{[P1$\langle$GM]}}
                \texttt{OUTPUT: False} \\
            }
        }
    }
     \\ \\

    \theutterance \stepcounter{utterance}  
    & \multicolumn{4}{p{0.6\linewidth}}{
        \cellcolor[rgb]{0.561,0.737,0.561}{
            \makecell[{{p{\linewidth}}}]{
                \texttt{\tiny{[P1$\rangle$GM]}}
                \texttt{\{"input": 10, "predicted\_output": 0, "current\_mode": "probing", "input\_rationale": "Testing with 10. Previous results: f(0)=0, f(1)!=1, f(2)!=4, f(3)!=9, f(4)!=16, f({-}1)!=1, f({-}2)!=0, f(5)!=0. The function is not identity, square, cube, or constant zero. It seems to return 0 for 0, but non{-}zero for others? Or maybe it returns something specific like the number of set bits, or a hash. Let's try a larger number to see if there's a pattern with magnitude.", "hypothesis": "unknown", "confidence\_score": 10\}} \\
            }
        }
    }
    & & \\ \\

    \theutterance \stepcounter{utterance}  
    & & & \multicolumn{4}{p{0.6\linewidth}}{
        \cellcolor[rgb]{0.9,0.9,0.9}{
            \makecell[{{p{\linewidth}}}]{
                \texttt{\tiny{[P1$\langle$GM]}}
                \texttt{OUTPUT: False} \\
            }
        }
    }
     \\ \\

    \theutterance \stepcounter{utterance}  
    & \multicolumn{4}{p{0.6\linewidth}}{
        \cellcolor[rgb]{0.561,0.737,0.561}{
            \makecell[{{p{\linewidth}}}]{
                \texttt{\tiny{[P1$\rangle$GM]}}
                \texttt{\{"input": 1, "predicted\_output": 0, "current\_mode": "probing", "input\_rationale": "Testing with 1 again, but predicting 0. Previous f(1)!=1. If f(1)=0, it might be a function that returns 0 for all positive integers, or something similar. Let's check if f(1) is 0.", "hypothesis": "unknown", "confidence\_score": 10\}} \\
            }
        }
    }
    & & \\ \\

    \theutterance \stepcounter{utterance}  
    & & & \multicolumn{4}{p{0.6\linewidth}}{
        \cellcolor[rgb]{0.9,0.9,0.9}{
            \makecell[{{p{\linewidth}}}]{
                \texttt{\tiny{[P1$\langle$GM]}}
                \texttt{OUTPUT: False} \\
            }
        }
    }
     \\ \\

    \theutterance \stepcounter{utterance}  
    & \multicolumn{4}{p{0.6\linewidth}}{
        \cellcolor[rgb]{0.561,0.737,0.561}{
            \makecell[{{p{\linewidth}}}]{
                \texttt{\tiny{[P1$\rangle$GM]}}
                \texttt{\{"input": 1, "predicted\_output": 2, "current\_mode": "probing", "input\_rationale": "Testing with 1 again, predicting 2. Previous f(1)!=1. If f(1)=2, it might be x+1. Let's check if f(x) = x+1. f(0)=1? But f(0)=0 was True. So not x+1. Maybe it's related to bit manipulation. Let's try a different approach. What if f(x) is the number of 1s in binary representation? f(0)=0 (True). f(1)=1 (False, so not popcount). What if f(x) is x XOR 1? f(0)=1 (False, f(0)=0). What if f(x) is 0 for even, 1 for odd? f(0)=0 (True). f(1)=1 (False). So not that. Let's try f(1)=2 again, just to see.", "hypothesis": "unknown", "confidence\_score": 10\}} \\
            }
        }
    }
    & & \\ \\

    \theutterance \stepcounter{utterance}  
    & & & \multicolumn{4}{p{0.6\linewidth}}{
        \cellcolor[rgb]{0.9,0.9,0.9}{
            \makecell[{{p{\linewidth}}}]{
                \texttt{\tiny{[P1$\langle$GM]}}
                \texttt{OUTPUT: False} \\
            }
        }
    }
     \\ \\

    \theutterance \stepcounter{utterance}  
    & \multicolumn{4}{p{0.6\linewidth}}{
        \cellcolor[rgb]{0.561,0.737,0.561}{
            \makecell[{{p{\linewidth}}}]{
                \texttt{\tiny{[P1$\rangle$GM]}}
                \texttt{\{"input": 1, "predicted\_output": 1, "current\_mode": "probing", "input\_rationale": "I previously tested f(1)=1 and it was False. I need to find what f(1) actually is. Let's try a value that is very different from 1, like a large number, to see if there's a pattern with magnitude or specific properties. However, since I have limited turns, let's try to deduce from f(0)=0. If f(x) = 0 for all x, then f(1) should be 0. I tested f(1)=0 and it was False. So f(1) is not 0, 1, 2. Let's try f(1)=3.", "hypothesis": "unknown", "confidence\_score": 10\}} \\
            }
        }
    }
    & & \\ \\

    \theutterance \stepcounter{utterance}  
    & & & \multicolumn{4}{p{0.6\linewidth}}{
        \cellcolor[rgb]{0.9,0.9,0.9}{
            \makecell[{{p{\linewidth}}}]{
                \texttt{\tiny{[P1$\langle$GM]}}
                \texttt{OUTPUT: False} \\
            }
        }
    }
     \\ \\

    \theutterance \stepcounter{utterance}  
    & \multicolumn{4}{p{0.6\linewidth}}{
        \cellcolor[rgb]{0.561,0.737,0.561}{
            \makecell[{{p{\linewidth}}}]{
                \texttt{\tiny{[P1$\rangle$GM]}}
                \texttt{\{"input": 1, "predicted\_output": 3, "current\_mode": "probing", "input\_rationale": "Testing f(1)=3. Previous results: f(0)=0, f(1)!=1, f(1)!=0, f(1)!=2. If f(1)=3, it might be x+2? But f(0)=2, which contradicts f(0)=0. So not x+2. Maybe it's a specific constant for x=1? Or maybe f(x) = x\^{}2 + x? f(0)=0, f(1)=2 (False). f(x) = x\^{}3? f(1)=1 (False). Let's try f(1)=3 to see if it's a pattern like f(x) = x + 2 for x>0? But f(0)=0. This is confusing. Let's try a different input entirely. What if f(x) is related to the number of divisors? f(0) is undefined or 0. f(1)=1 (False). What if f(x) is the largest power of 2 less than or equal to x? f(0)=0, f(1)=1 (False). Let's try f(1)=3.", "hypothesis": "unknown", "confidence\_score": 10\}} \\
            }
        }
    }
    & & \\ \\

    \theutterance \stepcounter{utterance}  
    & & & \multicolumn{2}{p{0.3\linewidth}}{
        \cellcolor[rgb]{0.9,0.9,0.9}{
            \makecell[{{p{\linewidth}}}]{
                \texttt{\tiny{[GM$|$GM]}}
                \texttt{Game Verdict: LOSS (turn limit reached)} \\
            }
        }
    }
    & & \\ \\

\end{supertabular}
}

\end{document}
